\hmsubsection{System Architecture Overview}{FS}{}\label{sec:system-architecture-overview}
This section provides a high-level overview of the system. The detailed
hardware architecture is described in Section \ref{sec:hardware-architecture},
the software architecture in Section \ref{sec:software-architecture},
and the data flow in Section \ref{sec:data-flow}.
The system consists of three computing nodes and one electrical
subsystem. Each node has a distinct responsibility:
\begin{itemize}
 \item \textbf{OAK-D S2 camera.} Captures images of the workspace and
 transmits them to the Raspberry Pi over USB 3.0.
 \item \textbf{Raspberry Pi 5.} Runs the vision pipeline, the inverse
 kinematics solver, and the cycle supervisor. This node performs the
 high-level decision-making.
 \item \textbf{OpenRB-150 motor controller.} Receives goal positions
 from the Pi over USB serial and drives the five Dynamixel servos.
\end{itemize}
The electrical subsystem is the 12 V bench supply that powers the
motors. It is described in Section \ref{sec:power-distribution}.
The separation between the Pi and the OpenRB-150 is the central
architectural decision. The Pi runs Linux, which is suitable for vision
and calibration code but less appropriate for deterministic motor timing.
The OpenRB-150 runs bare-metal C++, which is suitable for tight motor
timing but not for running OpenCV. This division assigns each node to
the tasks for which it is best suited.
The boundary between the two nodes is the USB serial link. The Pi
sends short text commands and the controller replies with status
tokens. The protocol is described in Section \ref{sec:comm}. Text was
selected instead of a binary frame because it is easier to debug from a
serial terminal, and the cable length and bench environment did not
justify adding a CRC check in the current setup. This is one of the
differences from Sortify \cite{sortify}, which used CRC-16 because the
arm was intended for installation on a moving vehicle.
Chapter \ref{cha:implementation} describes each node in detail.
